% reality/part_intro-DE.tex

\chapter*{Realität}
\phantomsection
\addcontentsline{toc}{chapter}{Realität}

Dieser Teil legt dar, wie kanonischer Zustand und Laufzeitzustand mit der
gebauten und beobachteten Eisenbahnrealität verbunden werden.

Die vorangehenden Teile haben konstruktive Alignment-Semantik, metrische
Realisierung, Eisenbahnzustand und Laufzeitbewertung begründet. Der vorliegende
Teil erklärt, wie diese genehmigten Begriffe in der Ingenieurpraxis
eingebettet, beobachtet, realisiert und verglichen werden.

Realität definiert vorgelagerte Begriffe nicht neu. Sie verwendet sie und
begründet die Abfolge
\[
\begin{aligned}
&\text{konstruktive Alignment-Semantik}
\rightarrow
\text{metrische Realisierung}
\\
&\rightarrow
\text{World-Einbettung}
\\
&\rightarrow
\text{physische Realisierung}
\\
&\rightarrow
\text{Messung und Beobachtung}
\rightarrow
\substack{\text{Ziel--Realisiert-}\\\text{Vergleich}}
\rightarrow
\text{ingenieurtechnische Interpretation}.
\end{aligned}
\]

Diese Abfolge wahrt die Unterscheidung zwischen konstruktiver Identität,
metrischer Realisierung, CRS-Kontext, physischer Realisierung, realisiertem
Eisenbahnzustand, Messung, Darstellung und operativer Interpretation.

\section*{Bezug zur AIM-Roadmap}

In der AIM-Roadmap ist Realität die Ebene, auf der zur Laufzeit bewerteter
Zustand auf gemessene und gebaute Infrastruktur trifft.

Die Orientierungsfrage der folgenden Abbildung lautet: Welche Ebenenfolge
verhindert, dass Ziel-, realisierte und beobachtete Zustände ineinanderfallen?

\begin{figure}[ht]
\centering
% figures/reality/icon_reality_target_realized_observed-DE.tex
\begin{tikzpicture}[
  node distance=1.3cm,
  box/.style={draw, rounded corners, minimum width=3.9cm, minimum height=0.9cm, align=center},
  arrow/.style={->, thick}
]
\node[box] (target) {Ziel};
\node[box, below=of target] (realized) {Realisiert};
\node[box, below=of realized] (observed) {Beobachtet};
\draw[arrow] (target) -- (realized);
\draw[arrow] (realized) -- (observed);
\draw[thin] ($(realized.center)+(0,0)$) ellipse (2.3cm and 0.62cm);
\draw[dashed, thin] ($(observed.center)+(0,0)$) ellipse (2.6cm and 0.72cm);
\end{tikzpicture}

\caption{Realitätssymbol: Ziel-, realisierter und beobachteter Eisenbahnzustand als abhängige Ebenen.}
\label{fig:icon-reality-target-realized-observed}
\end{figure}

\noindent
Die folgenden Kapitel verwenden diese Trennung, um Einbettung, Messung und
Vergleich zu strukturieren.

\begin{principle}[Realitätsprinzip]
\label{princ:reality-principle}
In AIM sind Modellzustand und physischer Eisenbahnzustand miteinander
verbunden, aber begrifflich verschieden; sie werden durch Realisierungs-,
Mess- und Interpretationsoperatoren verknüpft.
\end{principle}

\begin{summary}
Realität legt dar, wie genehmigte Alignment-Semantik und Laufzeitzustand in
die ingenieurtechnische Welt eingebettet, durch Daten beobachtet und durch den
Vergleich von Ziel und Realisiert interpretiert werden.
\end{summary}

Diese Trennung bereitet den Teil Modellierung vor, in dem persistente
Strukturen so entworfen werden, dass sie genau diese Unterscheidungen in
operativen Informationssystemen bewahren.
